\chapter{Junk Food and Your Health}

\section*{Definition and Overview}
Junk food is a term for foods and drinks that are high in kilojoules, sugar, salt, and unhealthy fats, and low in nutrients, fibre, and other beneficial components. Examples include fast food, chips, confectionery, soft drinks, cakes, biscuits, and many processed snacks. These foods are convenient, cheap, and heavily marketed, and they make up a large part of many Australian diets. Regular consumption contributes to weight gain, obesity, heart disease, type 2 diabetes, and dental decay. The Australian Dietary Guidelines recommend limiting discretionary foods to small amounts, and reducing junk food is one of the most effective steps for improving health.

\section*{Epidemiology}
Australians obtain around a third of their energy from discretionary foods, which are the foods and drinks not needed for a healthy diet. This is well above the recommended limit of 10 per cent. Around two-thirds of Australian adults are overweight or obese, and one in four children, with junk food a major contributor. Soft drink consumption has fallen in recent years but remains high, and fast food is eaten by around a third of Australians each week. The health costs of poor diet are substantial, contributing to heart disease, diabetes, and cancer.

\section*{Aetiology and Pathophysiology}
Junk food affects health through several mechanisms. It is energy-dense but low in nutrients, so it is easy to overconsume kilojoules while missing vitamins, minerals, and fibre. Sugar, especially in drinks, is absorbed quickly and, in excess, promotes weight gain, insulin resistance, and type 2 diabetes, and it feeds the bacteria that cause dental decay. Unhealthy fats, including saturated and trans fats, raise cholesterol and heart disease risk, while salt raises blood pressure. Highly processed foods are designed to be palatable and are linked to overeating through effects on appetite and reward pathways. Over time, these effects compound to drive obesity and chronic disease.

\section*{Clinical Features}
\begin{itemize}
    \item Weight gain and obesity
    \item Raised blood pressure, cholesterol, and blood glucose
    \item Increased risk of heart disease, stroke, and type 2 diabetes
    \item Dental decay and gum disease
    \item Fatigue and reduced concentration from poor nutrition
    \item Digestive problems from low fibre
    \item A large proportion of daily energy coming from discretionary foods
    \item Effects on children: obesity, poor eating habits, and school performance
\end{itemize}

\section*{Differential Diagnosis}
The health effects of junk food are assessed against the overall diet.
\begin{itemize}
    \item \textbf{Energy-dense, nutrient-poor foods:} the definition of junk food
    \item \textbf{Sugar-sweetened drinks:} a major source of excess sugar
    \item \textbf{Fast food:} high in kilojoules, salt, and unhealthy fats
    \item \textbf{Highly processed foods:} engineered for overconsumption
    \item \textbf{Occasional versus regular consumption:} frequency determines impact
    \item \textbf{Other contributors to weight and health:} inactivity, genetics, and medical conditions
\end{itemize}

\section*{When to Seek Medical Care}
Most people do not need medical care to reduce junk food, but a doctor or dietitian can help with weight management, diabetes, high blood pressure, or cholesterol. Children's diets and weight should be reviewed at health checks, and anyone with symptoms of diabetes or other diet-related disease should be assessed. Dietitians provide individualised guidance for lasting change.

\textbf{Call 000 (or local emergency services) for:}
\begin{itemize}
    \item Chest pain, sudden weakness, or difficulty speaking, which may indicate heart attack or stroke
    \item Very high blood sugar with vomiting, abdominal pain, or confusion
    \item Difficulty breathing or collapse
    \item Signs of a serious reaction, including severe allergic reactions
\end{itemize}

\section*{Diagnosis and Investigations}
\begin{itemize}
    \item \textbf{Dietary assessment:} review of typical food and drink intake
    \item \textbf{Weight and body mass index:} tracking over time
    \item \textbf{Blood pressure:} screening for hypertension
    \item \textbf{Blood tests:} glucose, cholesterol, and liver function
    \item \textbf{Dental review:} for decay and gum disease
    \item \textbf{Screening in children:} growth checks and dietary review
    \item \textbf{Follow-up:} monitoring of weight and metabolic health
\end{itemize}

\section*{Management}
\begin{itemize}
    \item \textbf{Limiting discretionary foods:} aiming for less than 10 per cent of energy
    \item \textbf{Replacing sugary drinks:} with water, plain milk, or reduced options
    \item \textbf{Cooking more at home:} using whole foods
    \item \textbf{Planning snacks:} fruit, vegetables, nuts, and yoghurt instead of packaged snacks
    \item \textbf{Reading labels:} checking kilojoules, sugar, salt, and saturated fat
    \item \textbf{Reducing portion sizes:} and not supersizing
    \item \textbf{Family meals:} modelling healthy eating for children
    \item \textbf{Professional support:} dietitians and weight management programs
    \item \textbf{Treating consequences:} management of weight, diabetes, and dental disease
\end{itemize}

\section*{Complications}
\begin{itemize}
    \item Obesity and its many effects
    \item Type 2 diabetes and insulin resistance
    \item Heart disease, stroke, and high blood pressure
    \item Dental decay and tooth loss
    \item Fatty liver disease
    \item Some cancers linked to diet
    \item Effects on mental health, energy, and sleep
    \item Costs to individuals and the health system
\end{itemize}

\section*{Prognosis and Outcomes}
Reducing junk food improves health quickly and progressively. Blood pressure, cholesterol, and glucose improve with weight loss, dental decay is preventable, and the risk of heart disease and diabetes falls substantially. Most people find that small, consistent changes, such as cutting sugary drinks and improving snacks, produce lasting benefits, and these changes are most effective when made as a family.

\section*{Prevention and Self-Care}
\begin{itemize}
    \item Drink water instead of sugary drinks
    \item Keep healthy snacks available and junk food out of sight
    \item Eat regular meals, and avoid eating in front of screens
    \item Cook more meals at home from whole foods
    \item Read food labels and choose lower-sugar, lower-salt options
    \item Limit fast food to occasional treats
    \item Serve sensible portions and avoid second helpings
    \item Be a healthy role model for children
    \item Plan ahead for work, school, and travel snacks
    \item Seek dietetic advice if you are struggling with weight or health
\end{itemize}

\section*{Sources and Further Reading}
The following reputable sources provide further information for healthcare students and patients seeking reliable, current advice about junk food and healthy eating.
\begin{itemize}
    \item Eat For Health (NHMRC). \textit{Australian dietary guidelines.} https://www.eatforhealth.gov.au/
    \item Dietitians Australia. \textit{Healthy eating and discretionary foods.} https://dietitiansaustralia.org.au/
    \item Healthdirect Australia. \textit{Healthy eating and weight.} https://www.healthdirect.gov.au/
    \item LiveLighter. \textit{Campaigns to reduce junk food and sugary drinks.} https://livelighter.com.au/
    \item Heart Foundation Australia. \textit{Heart-healthy eating.} https://www.heartfoundation.org.au/
    \item World Health Organization. \textit{Healthy diet.} https://www.who.int/
\end{itemize}
